% GENERATED FILE. Edit canonical lessons, then run npm run generate:books.
% canonical-source-hash: fnv1a64:a80c685c746b0bce
% canonical-lessons: GU-C116-tarbuc, GU-C116-papaiyu, GU-C116-jamphal, GU-C116-magphali, GU-C116-dhokla

\chapter{Watermelon, Papaya, Guava, Peanuts, Dhokla}
\label{ch:gu-watermelon-papaya-guava-peanuts-dhokla}
\hlchaptermodality{116}

\begin{chapteropening}
\textbf{By the end of this chapter:} \emph{I can say a watermelon, a papaya, a guava, peanuts and dhokla.}

Complete the last lesson of chapter 116: i can say a watermelon, a papaya, a guava, peanuts and dhokla.
\end{chapteropening}

\section[tarbūc]{\gu{તરબૂચ} (tarbūc) — a watermelon}

\label{lesson:GU-C116-tarbuc}



\begin{quote}
Before the new one: say the Gujarati for a lemon, then the Gujarati for a coconut.
\end{quote}

\subsection*{You'll want to know: \gu{તરબૂચ}}
\textbf{\gu{તરબૂચ}} (\emph{tarbūc}) — \textquotedblleft{}a watermelon\textquotedblright{}.

\begin{grammarlens}[title={What you've built}]
One more word for this chapter.
\end{grammarlens}

\subsection*{Your turn}
\begin{itemize}
\raggedright
  \item \emph{Say it:} \emph{tarbūc}
  \item \emph{Say it:} \emph{tarbūc}, once more
  \item \emph{Say it:} say \emph{tarbūc}
  \item \emph{Recall:} say the Gujarati for a lemon, then the Gujarati for a coconut, then say \emph{tarbūc} again
\end{itemize}

\begin{tcolorbox}[breakable,colback=teal!4,colframe=teal!35!black,title={Before you move on}]
What does \gu{તરબૂચ} mean? (\textquotedblleft{}A watermelon\textquotedblright{}.) Say it once more.
\end{tcolorbox}
\section[papaiyũ]{\gu{પપૈયું} (papaiyũ) — a papaya}

\label{lesson:GU-C116-papaiyu}



\begin{quote}
Before the new one: say the Gujarati for a coconut, then the Gujarati for a watermelon.
\end{quote}

\subsection*{You'll want to know: \gu{પપૈયું}}
\textbf{\gu{પપૈયું}} (\emph{papaiyũ}) — \textquotedblleft{}a papaya\textquotedblright{}.

\begin{grammarlens}[title={What you've built}]
One more word for this chapter.
\end{grammarlens}

\subsection*{Your turn}
\begin{itemize}
\raggedright
  \item \emph{Say it:} \emph{papaiyũ}
  \item \emph{Say it:} \emph{papaiyũ}, once more
  \item \emph{Say it:} say \emph{papaiyũ}
  \item \emph{Recall:} say the Gujarati for a coconut, then the Gujarati for a watermelon, then say \emph{papaiyũ} again
\end{itemize}

\begin{tcolorbox}[breakable,colback=teal!4,colframe=teal!35!black,title={Before you move on}]
What does \gu{પપૈયું} mean? (\textquotedblleft{}A papaya\textquotedblright{}.) Say it once more.
\end{tcolorbox}
\section[jāmphaḷ]{\gu{જામફળ} (jāmphaḷ) — a guava}

\label{lesson:GU-C116-jamphal}



\begin{quote}
Before the new one: say the Gujarati for a watermelon, then the Gujarati for a papaya.
\end{quote}

\subsection*{You'll want to know: \gu{જામફળ}}
\textbf{\gu{જામફળ}} (\emph{jāmphaḷ}) — \textquotedblleft{}a guava\textquotedblright{}.

\begin{grammarlens}[title={What you've built}]
One more word for this chapter.
\end{grammarlens}

\subsection*{Your turn}
\begin{itemize}
\raggedright
  \item \emph{Say it:} \emph{jāmphaḷ}
  \item \emph{Say it:} \emph{jāmphaḷ}, once more
  \item \emph{Say it:} say \emph{jāmphaḷ}
  \item \emph{Recall:} say the Gujarati for a watermelon, then the Gujarati for a papaya, then say \emph{jāmphaḷ} again
\end{itemize}

\begin{tcolorbox}[breakable,colback=teal!4,colframe=teal!35!black,title={Before you move on}]
What does \gu{જામફળ} mean? (\textquotedblleft{}A guava\textquotedblright{}.) Say it once more.
\end{tcolorbox}
\section[magphaḷī]{\gu{મગફળી} (magphaḷī) — peanuts}

\label{lesson:GU-C116-magphali}



\begin{quote}
Before the new one: say the Gujarati for a papaya, then the Gujarati for a guava.
\end{quote}

\subsection*{You'll want to know: \gu{મગફળી}}
\textbf{\gu{મગફળી}} (\emph{magphaḷī}) — \textquotedblleft{}peanuts\textquotedblright{}.

\begin{grammarlens}[title={What you've built}]
One more word for this chapter.
\end{grammarlens}

\subsection*{Your turn}
\begin{itemize}
\raggedright
  \item \emph{Say it:} \emph{magphaḷī}
  \item \emph{Say it:} \emph{magphaḷī}, once more
  \item \emph{Say it:} say \emph{magphaḷī}
  \item \emph{Recall:} say the Gujarati for a papaya, then the Gujarati for a guava, then say \emph{magphaḷī} again
\end{itemize}

\begin{tcolorbox}[breakable,colback=teal!4,colframe=teal!35!black,title={Before you move on}]
What does \gu{મગફળી} mean? (\textquotedblleft{}Peanuts\textquotedblright{}.) Say it once more.
\end{tcolorbox}
\section[ḍhokḷā̃]{\gu{ઢોકળાં} (ḍhokḷā̃) — dhokla}

\label{lesson:GU-C116-dhokla}



\begin{quote}
Before the new one: say the Gujarati for a guava, then the Gujarati for peanuts.
\end{quote}

\subsection*{You'll want to know: \gu{ઢોકળાં}}
\textbf{\gu{ઢોકળાં}} (\emph{ḍhokḷā̃}) — \textquotedblleft{}dhokla\textquotedblright{}.

\begin{grammarlens}[title={What you've built}]
One more word for this chapter.
\end{grammarlens}

\subsection*{Your turn}
\begin{itemize}
\raggedright
  \item \emph{Say it:} \emph{ḍhokḷā̃}
  \item \emph{Say it:} \emph{ḍhokḷā̃}, once more
  \item \emph{Say it:} say \emph{ḍhokḷā̃}
  \item \emph{Recall:} say the Gujarati for a guava, then the Gujarati for peanuts, then say \emph{ḍhokḷā̃} again
\end{itemize}

\begin{tcolorbox}[breakable,colback=teal!4,colframe=teal!35!black,title={Before you move on}]
What does \gu{ઢોકળાં} mean? (\textquotedblleft{}Dhokla\textquotedblright{}.) Say it once more.
\end{tcolorbox}
